\section{波函数与薛定谔方程}
\label{sec:波函数与薛定谔方程}
波尔理论在解释氢原子光谱上是成功的，但是，它也存在诸多局限性，其实质上也只能解释氢原子的性质，而无法适用于结构更为复杂的原子。根本原因在于波尔理论只能算是一种半经典半量子的理论，其量子性在于它的定态假设给出的能级和分立轨道，其经典性在于它仍然保留经典理论中的电子绕核做圆周运动的想法，因此，其注定需要被更完备的量子理论所取代。

波函数，尽管听上去是一些很复杂的东西，但它的想法并没有那么复杂。德布罗意的物质波指出，任何实物粒子都同时具有波动特性，并通过\fancyref{fml:德布罗意关系式}给出了粒子和波动间的换算关系，但问题在于，德布罗意波到底是什么？即，德布罗意波中到底是什么物理量在发生波动？作为参考，在机械振荡中是位移在波动，在电磁振荡中则是电磁矢量在波动。


但总要承认的是，德布罗意波中总也得有什么东西在波动，我们称为波函数$\Psi(\vb*{r},t)$。

我们知道，平面简谐波的表达式是
\begin{Equation}
    y(x,t)=A\cos\qty[\omega t-kx]=A\cos\qty[2\pi\qty(\nu t-\frac{1}{\lambda}x)]
\end{Equation}
我们不妨将上式用欧拉公式改写为复振动形式
\begin{Equation}
    y(x,t)=A\exp[2\pi\i\qty(\nu t-\frac{1}{\lambda}x)]
\end{Equation}
波函数$\Psi(x,t)$既然也是一种波动，那么也应具有以上形式，当然我们会说上式似乎少了取实部的过程，这其实是不必要的，因为对于波函数$\Psi(x,t)$而言，我们就需要它表现为复振动。

波函数，因此，应当具有以下形式
\begin{Equation}
    \Psi(x,t)=\Psi_0\exp[2\pi\i\qty(\nu t-\frac{1}{\lambda}x)]
\end{Equation}
根据\fancyref{fml:德布罗意关系式}
\begin{Equation}
    \Psi(x,t)=\Psi_0\exp[2\pi\i\qty(\frac{E}{h}t-\frac{p}{h}x)]
\end{Equation}
我们考虑将$1/h$提出并与$2\pi$构成$2\pi/h=1/\hbar$
\begin{Equation}
    \Psi(x,t)=\Psi_0\exp[\i\hbar\qty(Et-px)]
\end{Equation}
而一般情况下，波函数$\Psi$是对于三维空间而言的$\Psi(\vb*{r},t)$，这只需要作一些小的改写。
\begin{OldBoxFormula}[波函数的形式]
    微观粒子的波函数具有以下形式
    \begin{Equation}
        \Psi(\vb*{r},t)=\Psi_0\exp[\i\hbar\qty(Et-\vb*{p}\cdot\vb*{r})]
    \end{Equation}
\end{OldBoxFormula}
现在我们仍然没有解释波函数$\Psi(\vb*{r},t)$的物理意义到底是什么，事实上，波函数$\Psi(\vb*{r},t)$本身的意义并不明确，但不同的是，波函数的模方$\abs{\Psi(\vb*{r},t)}^2$具有清楚的意义，它表示粒子在某一时刻$t$出现在某一空间位置$\vb*{r}=(x,y,z)$的概率密度，这就意味着它应当满足归一化条件
\begin{Equation}
    \Itnt[\R^3]\abs{\Psi(\vb*{r},t)}^2\dd{V}=1
\end{Equation}
波函数的统计诠释可以为我们理解微观粒子的波粒二象性提供了一种更清晰的图像
\begin{itemize}
    \item 经典概念中的粒子，其运动规律遵循牛顿运动定律并具有确定的运动轨道。
    \item 经典概念中的波，其实体弥散在整个空间中并存在某个在空间波动的物理量。
\end{itemize}
波函数的模方表征概率密度，这就是说德布罗意波实际是一种概率波。这就是说，粒子还是经典意义下的那个粒子，但是，粒子不再有确定的位置和运动轨道了，粒子是以一种在空间不同位置具有不同出现概率的方式弥散在空间中，从而形成波。但由于这里的波动量是概率，因此概率波的弥散也并不要求破坏粒子的整体性。这样，我们就可以调和波粒二象性的冲突了。

波函数在某种意义上是微观粒子的“位置函数”，只不过它给出的是粒子出现在各位置的概率密度而不是确切的位置，经典意义下位置函数对应的动力学方程是牛顿第二定律，那么，量子意义下波函数对应的波动方程是什么？这将由著名的薛定谔方程给出，这里直接给出其形式。
\begin{OldBoxEquation}[薛定谔方程]
    微观粒子的波函数是\uwave{薛定谔方程}的解
    \begin{Equation}
        \i\hbar\pdv{\Psi}{t}+\frac{\hbar^2}{2m}\laplacian\Psi-V(\vb*{r},t)\Psi=0
    \end{Equation}
    其中$m$表示微观粒子的质量，而$V(\vb*{r},t)$表示微观粒子的势能函数。
\end{OldBoxEquation}
有趣的是，薛定谔方程其实是输运方程，但或许是因为$\pdv*{\Psi}{t}$项前的系数$\i\hbar$为虚数，薛定谔方程仍然能给出波动的结果。总之，薛定谔方程描述了粒子状态随时间和空间变化所遵循的普遍规律，就像牛顿第二定律与之经典力学那样，薛定谔方程就是量子力学的基本方程。